\section{Background}
This section serves as background information, not directly related to the subject, but nonetheless important for understanding it. I aim to highlight some of the more general ideas and problems that are relevant to the described topics, so that this can be kept in mind during the thesis.
\subsection{Topology in Condensed Matter}

In quantum mechanics, we describe systems using \textit{quantum numbers}, numbers that characterize the state of a system. Usually, in \textit{few-body} quantum mechanics, we look at symmetries of our system, and the conserved quantities corresponding to them. The generators of these symmetries are related to operators by canonical quantization, and we describe the state as a sum of eigenstates of the symmetry generators. To exemplify this, look at the the hydrogen atom. In this case, the system has three separate symmetry generators, related to rotations around three perpendicular axes. Each of these correspond to a specific quantum number used to describe the state.

However, in \textit{many-body} quantum mechanics, we often do not have the luxury of symmetries. Of course, when treating a crystal as a perfectly periodic background of ions, we have a discrete translational symmetry, giving rise to Bloch bands, which gives us a quantum number. However, when we introduce defects into this lattice, we can no longer rely on this description either. It's an unfortunate truth that nature is messy. Once we attempt to scale up our models to contain a macroscopic amount of particles, the symmetries that we relied on for our quantum mechanical description tend to break. This is one of the reasons why many-body quantum physics is so difficult. 

Topological quantum numbers don't suffer from this problem. They are quantum numbers quantized on purely \textit{topological} grounds (as opposed to \textit{geometrical} grounds). This sounds vague at best, but should be clearer after some examples.

The simplest form of a topological quantum number is the winding number. We obviously require single-valuedness for the order parameter, and for systems in which the order parameter is complex-valued, this means that along a closed contour the phase angle may only change by an integer value of $2\pi$. This integer is the winding number, and may, by Stokes' theorem, be related to physically measurable quantities, which are referred to as a defect, because they require our order parameter to vanish. An example of a system with complex order parameter is a superfluid, in which the corresponding quantized quantity is circulation within the fluid. Another example is a superconductor, within which the magnetic flux is quantized. We say such quantities are \textit{topologically protected}, in the sense that we can't change the quantity by continuously changing the order parameter, as the winding number has to remain an integer. 

A slightly more complex example is the Quantum Hall Effect (QHE), in which the conductance is quantized. The topology of this system is more complicated, but again we can define a topological invariant, the Chern number, which corresponds to the amount of conductance quanta.

\begin{tcolorbox}[title = Sources, colback = gray!5, colframe = gray]
    \href{file:/Users/jobenning/Library/Mobile\ Documents/com\~apple\~CloudDocs/Desktop/Books\ \&\ Papers/Papers/Thouless\ -\ Introduction\ to\ Topological\ Quantum\ Numbers.pdf}{Thouless - Introduction to Topological Quantum Numbers}
\end{tcolorbox}

\subsection{Kitaev Chain}
The Kitaev chain Hamiltonian is given by
\begin{equation}
    \hat{H} = -t\sum_{j=1}^{N-1}(\hat{a}^\dagger_{j+1}\hat{a}_j+\hat{a}_j^\dagger \hat{a}_{j+1}) - \mu \sum_{j=1}^N\hat{a}^\dagger_j \hat{a}_j + \sum_{j=1}^{N-1}(\Delta^* \hat{a}_{j+1}\hat{a}_j + \Delta a_j^\dagger a_{j+1}^\dagger).
\end{equation}

It's important to understand that the pairing field $\Delta$ is not a self-consistently determined parameter, but rather taken to be a free parameter of the theory. This means that the model assumes that the chain is coupled to a superconductor, which serves as an electron reservoir, which allows the anomolous terms to exist. The anomolous term couples the holes and electrons, and is what causes the superconductivity. Due to a gauge symmetry in the definition of the fields we may take $\Delta$ real. 

This model can't support a steady current, so we would like to make the connection to the reservoir explicit, so that we can, by tuning the chemical potential of the different connection points, make current flow. 

\begin{tcolorbox}[title = Questions, colback = gray!5, colframe = black]
    \begin{enumerate}
        \item When do we use the basic model, and when do we use the model that attaches to the leads? \textcolor{blue}{We might have to connect the model to the leads if we want to find the current, because if we don't, the tip and the small Kitaev chain will just equilibrate, which makes the current zero.}
    \end{enumerate}
\end{tcolorbox}

 \begin{tcolorbox}[title = Sources, colback = gray!5, colframe = gray]
    \begin{enumerate} 
        \itemsep0em
        \item Doornenbal - Conductance of a Finite Kitaev Chain
        \item Kitaev - Unpaired Majorana fermions in quantum wires
        \item Leijnse - Introduction to topological superconductivity and Majorana fermions
    \end{enumerate}
    \end{tcolorbox}

\subsection{Keldysh Formalism}
\begin{tcolorbox}[title = Sources, colback = gray!5, colframe = gray]
    \begin{enumerate} 
        \itemsep0em
        \item Kamenev - Non-equilibrium Field Theory
        \item Stoof - Coherent Versus Incoherent Dynamics During Bose-Einstein Condensation in Atomic Gases
    \end{enumerate}
\end{tcolorbox}

